\chapter{Fundamentals}
\label{chap:fundamentals}

This chapter presents the theoretical knowledge that is needed for the
solution of the problem. Only contents that are used later in the report
are included. Section~2.1 revisits feed forward neural networks and their
training. Section~2.2 formalises continual learning and its standard
scenarios. Section~2.3 explains catastrophic forgetting in more detail, and
Section~2.4 defines the metrics that are used to measure it. Section~2.5
introduces the Mixture of Experts architecture and its gating mechanism.

\section{Feed Forward Neural Networks}
\label{sec:fnn}

An artificial neural network is a parametric function
$f_{\theta} : \mathbb{R}^{d} \rightarrow \mathbb{R}^{c}$ that is built from
layers of simple computational units. A fully connected layer computes
\begin{equation}
  h^{(l)} = \phi\!\left(W^{(l)} h^{(l-1)} + b^{(l)}\right),
  \label{eq:layer}
\end{equation}
where $h^{(l-1)} \in \mathbb{R}^{d_{l-1}}$ is the input of the layer,
$W^{(l)} \in \mathbb{R}^{d_l \times d_{l-1}}$ is the weight matrix,
$b^{(l)} \in \mathbb{R}^{d_l}$ is the bias vector, and $\phi$ is a non
linear activation function. Stacking several such layers gives a multi
layer perceptron (MLP). In this project, the rectified linear unit
$\phi(z) = \max(0, z)$, short ReLU, is used for the hidden layers, because
it is simple, fast and reduces the vanishing gradient problem
\cite{goodfellow2016}.

For a classification problem with $c$ classes, the last layer produces a
vector of logits $z \in \mathbb{R}^{c}$, which is converted into a
probability distribution by the softmax function,
\begin{equation}
  p(y = i \mid x) = \frac{\exp(z_i)}{\sum_{j=1}^{c} \exp(z_j)}.
  \label{eq:softmax}
\end{equation}
The network is trained by minimising the cross entropy loss
\begin{equation}
  \mathcal{L}_{\mathrm{CE}}(\theta)
    = -\frac{1}{N} \sum_{n=1}^{N} \log p(y_n \mid x_n; \theta)
  \label{eq:ce}
\end{equation}
over a training set $\{(x_n, y_n)\}_{n=1}^{N}$ with mini batch stochastic
gradient descent or one of its adaptive variants. In this project, the Adam
optimiser \cite{kingma2015} is used. Adam keeps an estimate of the first
and the second moment of the gradient for every parameter and adapts the
effective step size accordingly, which makes the training robust against a
suboptimal choice of the learning rate.

An important property of such networks for this project is that the
knowledge of the network is stored in a distributed way. Every weight
takes part in the computation of every output, and every output depends on
many weights. This property is one of the main reasons for the forgetting
problem that is described in Section~2.3.

\section{Continual Learning}
\label{sec:cl}

\subsection{Problem Definition}

In continual learning, the training data is not available as one static
data set but arrives as a sequence of $T$ tasks
$\mathcal{T}_1, \mathcal{T}_2, \dots, \mathcal{T}_T$. Each task
$\mathcal{T}_t$ consists of a data set
$\mathcal{D}_t = \{(x_i^{(t)}, y_i^{(t)})\}_{i=1}^{N_t}$
from a task specific distribution. At time $t$, the learner has only access
to $\mathcal{D}_t$, or at most to a very limited memory of earlier tasks,
and has to update its parameters from $\theta_{t-1}$ to $\theta_t$. The
goal is to perform well on all tasks seen so far, not only on the current
one.

This setting is fundamentally different from multi task learning, where the
data of all tasks is available at the same time and can be mixed freely,
and also from transfer learning, where only the performance on a new target
task matters and the performance on the source task is allowed to degrade.

\subsection{The Three Continual Learning Scenarios}
\label{sec:scenarios}

Van de Ven and Tolias \cite{vandeven2019} distinguish three scenarios of
increasing difficulty, which differ in the information that is available at
test time:

\begin{description}[style=nextline]
  \item[Task incremental learning.] The task identity is known at test
        time. The model may use task specific components, for example a
        separate output head per task, and it only has to solve the sub
        problem of the current task. SplitMNIST with one output head per
        task belongs to this scenario.
  \item[Domain incremental learning.] The task identity is not needed at
        test time, but all tasks share the same output space. The model
        must solve every task without knowing which task an input belongs
        to, but it never has to distinguish classes of different tasks.
        PermutedMNIST with one shared ten class head is the standard
        example.
  \item[Class incremental learning.] The model additionally has to infer
        the task identity itself, that is, it has to distinguish all
        classes of all tasks at the same time. This is the hardest of the
        three scenarios.
\end{description}

The two benchmarks of this project cover the first two scenarios. This is a
deliberate choice, as explained in Section~1.2: the main goal is to study
the stability plasticity behaviour of the Mixture of Experts architecture
in a clean setting, not to solve the harder class incremental problem,
which usually needs additional memory based techniques.

\subsection{The Stability Plasticity Dilemma}
\label{sec:dilemma}

Every continual learning method has to balance plasticity, the ability to
gain new knowledge, against stability, the ability to keep old knowledge
\cite{mermillod2013}. Regularisation based methods implement this trade off
explicitly with an additional stability term in the loss function.
Architectural methods implement it implicitly by assigning different parts
of the network to different tasks. The Mixture of Experts approach
investigated here realises the trade off through routing: the plasticity is
concentrated in the experts that the router selects for the new task, while
the remaining experts act as stable memory.

\section{Catastrophic Forgetting}
\label{sec:forgetting}

Catastrophic forgetting denotes the fast and severe loss of previously
learned performance when a neural network is trained on new information.
The phenomenon was documented systematically by McCloskey and Cohen
\cite{mccloskey1989}, who showed that networks trained sequentially on two
lists of associations perform almost at chance level on the first list
after learning the second one. French \cite{french1999} related the effect
to the distributed representations of connectionist models: because all
weights take part in the storage of every pattern, writing a new pattern
into the network inevitably overwrites parts of the old ones.

From the optimisation point of view, the reason is simple. During the
training on task $\mathcal{T}_t$, gradient descent follows the loss
landscape of $\mathcal{L}_t$ only. Nothing in the objective prevents the
parameters from moving to a region of the weight space that is good for
$\mathcal{T}_t$ but arbitrarily bad for the earlier tasks. Because the loss
surfaces of different tasks usually do not share a common minimum in a high
dimensional weight space, sequential optimisation almost guarantees
interference \cite{goodfellow2013}.

It is important to note that catastrophic forgetting is not caused by a
lack of capacity. Even networks with far more parameters than needed for
each single task forget, because the same parameters are reused for all
tasks. This observation is the key motivation for the present project: a
Mixture of Experts model can partition its capacity so that different tasks
rely on different parameters, and can therefore attack the root cause of
the interference instead of only dampening its symptoms.

\section{Metrics for Continual Learning}
\label{sec:clmetrics}

To make forgetting measurable, the evaluation follows the standard protocol
of Lopez-Paz and Ranzato \cite{lopezpaz2017}. Let
\begin{equation}
  R_{i,j} \in [0,1]
\end{equation}
be the test accuracy on task $\mathcal{T}_j$ after the model has been
trained on task $\mathcal{T}_i$, with $j \leq i$. The triangular matrix $R$
describes the complete learning dynamic of a run. Three scalar metrics are
derived from it.

The average accuracy is the mean performance over all tasks after the last
training stage,
\begin{equation}
  A = \frac{1}{T} \sum_{j=1}^{T} R_{T,j}.
  \label{eq:avgacc}
\end{equation}
It reflects plasticity and stability together and is the primary metric of
a continual learner.

The average forgetting measures, for every task, the difference between the
best accuracy ever reached on it and its accuracy at the end of training:
\begin{equation}
  F = \frac{1}{T-1} \sum_{j=1}^{T-1}
      \Big( \max_{i \in \{j, \dots, T-1\}} R_{i,j} - R_{T,j} \Big).
  \label{eq:forgetting}
\end{equation}
A value of $F = 0$ means perfect retention; larger positive values mean
more severe forgetting.

The backward transfer (BWT) measures the average influence that learning a
new task has on the performance on previous tasks,
\begin{equation}
  \mathrm{BWT} = \frac{1}{T-1} \sum_{j=1}^{T-1} \left( R_{T,j} - R_{j,j} \right).
  \label{eq:bwt}
\end{equation}
Negative values indicate forgetting; positive values would mean that
learning later tasks even improved the earlier ones.

\section{Mixture of Experts}
\label{sec:moe}

\subsection{Historical Roots}

The Mixture of Experts idea goes back to Jacobs, Jordan, Nowlan and
Hinton \cite{jacobs1991}, who proposed to replace one single network by a
set of expert networks together with a gating network. Each expert
specialises on a region of the input space, and the gating network learns
to assign every input to the most suitable expert. Because the experts are
trained jointly but compete through the gate, the architecture shows a
form of divide and conquer: a complex problem is decomposed into simpler
sub problems that are solved by specialists. Hierarchical variants of the
model were analysed afterwards by Jordan and Jacobs \cite{jordan1994}.

\subsection{Sparsely Gated Mixture of Experts}
\label{sec:sparsemoe}

The modern, scalable form of the idea is the sparsely gated MoE layer of
Shazeer et al.\ \cite{shazeer2017}. Instead of weighting all experts, the
gate activates only the $k$ best ones, typically $k \in \{1, 2\}$, which
keeps the computational cost per input almost constant while the total
number of parameters grows with the number of experts. This architecture
later became a core component of very large language models, most notably
the Switch Transformer \cite{fedus2022}, and showed that expert based
capacity scaling works reliably in practice.

Formally, an MoE layer consists of $E$ expert networks $f_1, \dots, f_E$
and a gating network $g$. For an input $h$, the gate first computes logits
and converts them into a probability distribution over the experts:
\begin{equation}
  g(h) = \mathrm{softmax}\!\left(W_g h + b_g\right) \in \mathbb{R}^{E}.
  \label{eq:gate}
\end{equation}
Only the $k$ experts with the largest gate values are evaluated. If
$\mathcal{S}(h) = \operatorname{top}\text{-}k\big(g(h)\big)$ denotes the
index set of the selected experts and $\tilde{g}_e$ their renormalised gate
values, the output of the layer is the convex combination
\begin{equation}
  \mathrm{MoE}(h) = \sum_{e \in \mathcal{S}(h)} \tilde{g}_e(h) \, f_e(h),
  \qquad
  \tilde{g}_e(h) = \frac{g_e(h)}{\sum_{e' \in \mathcal{S}(h)} g_{e'}(h)}.
  \label{eq:moe}
\end{equation}
All parameters, those of the experts and those of the gate, are trained
jointly by backpropagation; the gate receives gradients through the weights
$\tilde{g}_e$.

\subsection{Load Balancing}
\label{sec:loadbalance}

A well known failure mode of MoE training is expert collapse: the gate
quickly prefers one or two experts, these experts receive more gradient
signal and become even better, and the remaining experts starve. To prevent
this, an auxiliary loss is used that rewards a uniform usage of the experts
\cite{shazeer2017}. The variant used in this project follows the Switch
Transformer \cite{fedus2022}. Let $f_e$ be the fraction of samples of a
batch whose highest ranked expert is $e$, and let $P_e$ be the mean gate
probability of expert $e$ over the batch. Then
\begin{equation}
  \mathcal{L}_{\mathrm{LB}} = E \cdot \sum_{e=1}^{E} f_e \, P_e
  \label{eq:lb}
\end{equation}
is minimised when all experts are used equally. It is added to the task
loss with a small coefficient $\alpha$:
\begin{equation}
  \mathcal{L} = \mathcal{L}_{\mathrm{CE}} + \alpha \, \mathcal{L}_{\mathrm{LB}}.
  \label{eq:totalloss}
\end{equation}
As Chapter~7 will show, this auxiliary loss has an unexpected and important
side effect in the continual setting.

\subsection{Position of the MoE Layer in the Network}
\label{sec:moeposition}

A design question that is rarely discussed in the MoE literature, but is
central for this project, is where in the network the MoE layer should be
placed. In large language models, the MoE layers replace the feed forward
blocks in the middle of a deep stack, so the router always sees a hidden
representation that was produced by shared layers in front of it. Two
extreme positions can be distinguished.

With \emph{hidden level routing}, the input first passes one or more shared
layers, and the router assigns the resulting hidden representation to the
experts. The shared layers form a common feature extractor for all inputs,
and the experts only specialise on the later processing stages. With
\emph{input level routing}, the router and the experts act directly on the
raw input, and only the mixed expert output passes shared layers.

The difference matters for continual learning, because routing can only
protect the parameters that lie behind the router. Parameters in front of
the router process every input of every task and are updated by every
task, so they behave exactly like the weights of a dense network. If two
tasks differ in the way the raw input must be interpreted, as for example
two different pixel permutations do, the decisive task specific knowledge
lives precisely in this early mapping. A router behind a shared trunk
cannot protect this mapping, no matter how well it separates the tasks in
the hidden space. This consideration is taken up again in the design
decisions of Chapter~4 and becomes directly relevant in the evaluation in
Chapter~7.

\subsection{Why MoE Is Attractive for Continual Learning}
\label{sec:whymoe}

The connection between MoE and continual learning is the central idea of
this project. Three properties make the architecture promising. First,
conditional computation: only a subset of the network is active for any
given input, so if the router separates tasks, the gradient updates on a
new task touch only a subset of the parameters, which directly limits
interference. Second, spare capacity: the total number of parameters scales
with the number of experts, not with the cost per input, so new tasks can
occupy capacity that was not needed by earlier tasks. Third, learned task
inference: the gate itself is a small classifier that learns, without being
told the task identity, to distinguish the input distributions of the
tasks. Whether these properties actually appear in a sequential training
regime, and how strongly they protect against forgetting, is exactly what
the experiments in Chapter~7 are designed to measure.
